% ---------------------------------------------------------------------------
%  preamble.tex — shared packages and macros for the paperopenmp2 document
% ---------------------------------------------------------------------------
\usepackage[utf8]{inputenc}
\usepackage[T1]{fontenc}
\usepackage{lmodern}
\usepackage{amsmath,amssymb}
% amsthm is incompatible with llncs (both define proof/theorem).
% Load it only when the fallback article class is active.
\makeatletter
\@ifclassloaded{llncs}{}{\RequirePackage{amsthm}}
\makeatother
\usepackage{booktabs}
\usepackage{xcolor}
\usepackage{graphicx}
\usepackage{caption}
\usepackage{subcaption}
\usepackage{algorithm}
\usepackage{algpseudocode}
\usepackage{listings}
\usepackage{microtype}
\usepackage{tikz}
\usepackage{pgfplots}
\pgfplotsset{compat=1.18}
\usetikzlibrary{%
  arrows.meta, positioning, calc, fit, backgrounds,
  shapes.geometric, decorations.pathreplacing,
  shadows.blur, matrix, patterns}
\usepackage[hidelinks]{hyperref}
\usepackage[nameinlink]{cleveref}
\usepackage{multirow}

% ---- colour palette (matches house style) ----------------------------------
\definecolor{cpu}{HTML}{2C6E9B}     % blue   -- CPU / host
\definecolor{gpu}{HTML}{B5651D}     % orange -- GPU / device
\definecolor{acc}{HTML}{3F7A3F}     % green  -- accents
\definecolor{soft}{HTML}{F1F4F8}    % light blue bg
\definecolor{softg}{HTML}{FBEFE3}   % light orange bg
\definecolor{codebg}{HTML}{F7F7F4}  % code background

\colorlet{cA}{cpu}
\colorlet{cB}{gpu}
\colorlet{cC}{acc}
\colorlet{cD}{cpu!60!black}
\definecolor{cGrey}{HTML}{7F8C8D}

% ---- tikz styles -----------------------------------------------------------
\tikzset{
  box/.style    = {rounded corners=2pt, draw=black!55, thick, align=center,
                   font=\small, inner sep=5pt, minimum height=8mm},
  cpubox/.style = {box, fill=soft,  draw=cpu!70},
  gpubox/.style = {box, fill=softg, draw=gpu!75},
  lane/.style   = {rounded corners=3pt, draw=black!25, thick, inner sep=7pt},
  flow/.style   = {-{Stealth[length=2.4mm]}, thick, draw=black!65},
  dflow/.style  = {-{Stealth[length=2.4mm]}, thick, draw=black!45, dashed},
  lbl/.style    = {font=\scriptsize\itshape, text=black!60, align=center},
  threadbox/.style = {rounded corners=2pt, draw=cA!60, thick, fill=soft,
                      align=center, font=\scriptsize, inner sep=4pt,
                      minimum width=18mm, minimum height=7mm},
  mergebox/.style  = {rounded corners=2pt, draw=cC!70, thick, fill=cC!8,
                      align=center, font=\small, inner sep=5pt, minimum height=8mm},
  statebox/.style  = {circle, draw=black!50, fill=white, thick,
                      font=\scriptsize, inner sep=2pt, minimum size=6mm},
  zonenode/.style  = {rounded corners=1pt, draw=black!40, fill=white,
                      font=\scriptsize\ttfamily, inner sep=3pt},
}

% ---- listings --------------------------------------------------------------
\lstset{
  basicstyle=\ttfamily\footnotesize,
  backgroundcolor=\color{codebg},
  breaklines=true, columns=fullflexible, keepspaces=true,
  frame=single, rulecolor=\color{black!20},
  framesep=4pt, xleftmargin=4pt,
  commentstyle=\color{acc},
  keywordstyle=\color{cpu}\bfseries,
  language=C++, showstringspaces=false,
  morekeywords={pragma,omp,parallel,for,schedule,dynamic,threadprivate}
}

% ---- theorem environments -------------------------------------------------
% llncs already provides theorem, lemma, proposition, definition, remark, proof.
% When the fallback article class is used, define them ourselves.
\makeatletter
\@ifclassloaded{llncs}{}{%
  \newtheorem{theorem}{Theorem}[section]%
  \newtheorem{lemma}[theorem]{Lemma}%
  \newtheorem{proposition}[theorem]{Proposition}%
  \newtheorem{definition}[theorem]{Definition}%
  \newtheorem{remark}{Remark}[section]%
}
\makeatother

% ---- macros ----------------------------------------------------------------
\newcommand{\tool}{\textsc{PtcTlmc}}
\newcommand{\ltsmin}{\textsc{LTSmin}}
\newcommand{\code}[1]{\texttt{\small #1}}
\newcommand{\clocks}{\mathcal{C}}
\newcommand{\Loc}{L}
\newcommand{\Init}{\ell_0}
\newcommand{\Inv}{\mathit{Inv}}
\newcommand{\Guard}{\mathit{Grd}}
\newcommand{\Zone}{\mathcal{Z}}
\newcommand{\DBM}{\mathit{DBM}}
\newcommand{\Frontier}{\mathcal{F}}
\newcommand{\Passed}{\mathcal{P}}
\newcommand{\Wait}{\mathcal{W}}
\newcommand{\extraLU}{\mathit{Extra}^+_{LU}}
\newcommand{\succfn}{\mathit{succ}}   % renamed to avoid clash with \succ (math symbol)
\newcommand{\subsumedby}{\preceq}
% ---- algorithmic extensions -----------------------------------------------
\algdef{SE}[PARFOR]{ParFor}{EndParFor}[1]{\textbf{parallel for} #1 \textbf{do}}{\textbf{end parallel for}}
